\chapter{Conclusiones}
\label{chap:conclusiones}

\lettrine{E}{ste} trabajo ha adaptado cinco NAS Parallel Benchmarks al modelo
de programación OMP4Py y los ha evaluado experimentalmente en Finisterrae~III
con intérpretes Python \emph{free-threaded}. La evaluación cubre los benchmarks
CG, EP, FT, IS y MG en las clases de problema S, W y A, con los cuatro modos de
ejecución cuando el coste lo permite y entre 1 y 32~hilos. La campaña
principal con Python~3.15.0b1t incluye 630~ejecuciones. La campaña comparativa
entre Python~3.13.13t, 3.14.4t y 3.15.0b1t añade 2340~ejecuciones, que cubren
tanto los modos interpretados como los compilados. Todas las ejecuciones finales
terminaron con verificación \texttt{SUCCESSFUL}.


\section{Hallazgos principales}
\label{sec:hallazgos}

\subsection{Cobertura funcional del port}
El port es funcional y cubre los caminos relevantes de OMP4Py. La clase~S se
ejecutó con los modos~0, 1, 2 y~3, validando los caminos
interpretados y compilados. La clase~W se evaluó con los modos compilados~2 y
3, y la clase~A con el modo~3. Esta estructura permite comprobar corrección sin
dedicar recursos a configuraciones que ya son varios órdenes de magnitud más
lentas en clases menores.

\subsection{El tipado Cython como factor dominante}
El tipado Cython es el factor de rendimiento dominante. En las clases S y W, el
modo~3 mejora al modo~2 en todos los benchmarks medidos, con factores que van
desde el orden de $15\times$ en FT hasta más de dos órdenes de magnitud en los
casos más favorables. Por tanto, para estos benchmarks el tipado no es una
optimización menor. Es la condición que convierte el código OMP4Py en una
ejecución competitiva.

\subsection{Escalabilidad de EP y FT}
La campaña principal obtiene speedups altos con Python~3.15.0b1t en los
benchmarks con estructura más favorable. En clase~A, EP y FT superan
$10\times$ de speedup, y FT alcanza su mejor punto con 16~hilos. Estos
resultados muestran que el port puede explotar paralelismo real cuando el kernel
tiene suficiente trabajo independiente. EP en modo~1 confirma además que la
eliminación del GIL produce paralelismo real al nivel del propio Python, aunque
esa ruta no sea competitiva en tiempo absoluto frente al modo tipado.

\subsection{Límites de CG, IS y MG}
CG, IS y MG no obtienen escalabilidad sostenida en modo~3. CG empeora al
aumentar hilos por el coste de reducciones, normas y barreras frecuentes. IS
solo mejora de forma moderada con pocos hilos y pierde eficiencia a 32~hilos.
MG combina accesos intensivos a memoria, niveles de malla de distinta
granularidad y sincronizaciones entre fases. En estos benchmarks, el tipado
mejora el tiempo absoluto, pero no basta para convertir más hilos en speedup.

\subsection{El tipado neutraliza la versión de Python}
En los modos compilados, Python~3.13t, 3.14t y 3.15t escalan casi igual, así que
la versión no cambia la conclusión práctica. En los modos interpretados el efecto
sí es grande. Python~3.13t escala mucho peor al aumentar los hilos en CG, FT, IS y
MG, hasta varias veces más lento a 32~hilos, y solo EP queda al margen por no
compartir apenas estado entre hilos. El tipado del modo~3 elimina esa penalización, porque saca el
bucle caliente del espacio de objetos Python. Donde la versión influye,
Python~3.15t es la más rápida. Como los modos compilados son los únicos con
tiempos prácticos, esta diferencia no cambia la conclusión principal del trabajo,
pero confirma la mejora de escalado de las builds \emph{free-threaded} posteriores
a la 3.13t.

\subsection{La afinidad del entorno como factor crítico}
La afinidad del entorno fue un factor experimental crítico, aunque OMP4Py no
implementa la fijación de hilos a núcleos. Las ejecuciones preliminares
heredaban variables de afinidad de OpenMP del entorno, y eso producía una
serialización aparente en Python~3.15t. El mecanismo, detallado en la
Sección~\ref{sec:profiling-afinidad}, es indirecto. NumPy carga
\texttt{libgomp}, que ante esas variables restringe la máscara de CPU del
proceso a un solo núcleo durante el \texttt{import}, y los hilos
\emph{free-threaded} heredan esa máscara. Al eliminar \texttt{OMP\_PROC\_BIND},
\texttt{OMP\_PLACES} y \texttt{GOMP\_CPU\_AFFINITY} por defecto en las campañas
finales, EP y FT pasaron a escalar correctamente. Esta corrección cambia la
interpretación del trabajo. El problema inicial no era una incapacidad de
Python~3.15t ni un fallo del port, sino una restricción de la planificación
impuesta desde fuera de la biblioteca.

\subsection{Escalabilidad real frente a paralelismo ineficiente}
El profiling distingue la escalabilidad real del paralelismo ineficiente. Las
mediciones con \texttt{perf stat} muestran que EP y FT usan varias CPUs y
reducen el tiempo. En FT, el IPC y los fallos de caché explican que el óptimo
aparezca en 16~hilos en lugar de 32. En CG y MG, el proceso también puede usar
varias CPUs, pero el IPC cae y el tiempo empeora. Por tanto, los malos speedups
no tienen una única causa. En unos casos el kernel escala, y en otros el
paralelismo existe pero la sincronización o la memoria dominan. El sistema de
muestreo de Python~3.15 cuantifica ese coste desde el nivel del intérprete. La
fracción de tiempo paralelo en espera de barreras y reducciones crece con el
número de hilos en los puntos medidos, y muestra que CG, FT, IS y MG crean y
destruyen equipos entre regiones paralelas, frente al equipo persistente de EP.
El trazado de Python~3.15 en Finisterrae~III confirma el gradiente estructural.
EP crea un único equipo, mientras que CG crea miles de hilos y decenas de miles
de entradas en barrera en una ejecución de clase~S.


\section{Limitaciones}
\label{sec:limitaciones}

Las amenazas a la validez se agrupan en la
Sección~\ref{sec:amenazas-validez} y se concretan después en la
Sección~\ref{sec:resultados-limitaciones}.
La primera limitación es experimental. La comparación entre versiones no incluye
clase~A, una decisión deliberada para contener el uso del cluster. La segunda es
de medida. La salida NAS redondea a dos decimales, de modo que algunos tiempos de
clase~S en IS y MG sirven para confirmar corrección y orden de magnitud, pero no
para ratios finos. La tercera afecta al profiling. \texttt{perf stat} mide el
proceso completo, y el profiler de Python~3.15 se usa como diagnóstico
estructural en modo~1 sobre clase~S. El trazado completa los cinco benchmarks y
el muestreo aporta los puntos válidos de la campaña. La cuarta es de alcance.
Los resultados caracterizan los cinco benchmarks adaptados, los nodos
\texttt{clk} y el rango de 1 a 32~hilos, pero no agotan otros tamaños,
arquitecturas ni políticas de afinidad.


\section{Trabajo futuro}
\label{sec:trabajo-futuro}

\subsection{Instrumentación y campañas complementarias}
La prioridad inmediata es cerrar la brecha entre diagnóstico y medición final.
El muestreo y el trazado de Python~3.15 ya identifican espera en barreras,
reducciones y creación de equipos, pero lo hacen desde fuera del runtime y sobre
modo~1. La siguiente campaña debería instrumentar OMP4Py por región paralela y en
los modos compilados, midiendo directamente tiempo útil, espera, reducciones y
creación de hilos. También convendría extender el muestreo a una matriz más fina
y a modos compilados, para ampliar la cobertura de plataforma, modo y
herramienta.

También queda una campaña selectiva de clase~A para versiones. No es necesario
repetir toda la matriz. Bastaría con EP y FT, modo~3, Python~3.13t y 3.14t, para
confirmar si los speedups de clase~W se mantienen con la carga mayor. Una segunda
campaña fina alrededor de los óptimos de FT e IS permitiría localizar mejor el
mejor número de hilos que el barrido por potencias de dos.

\subsection{Mejoras del runtime}
El candidato principal de mejora es reutilizar el equipo de hilos. El runtime
actual crea y destruye un equipo por región paralela, y en kernels con muchas
regiones cortas, como CG o MG, ese coste aparece miles de veces. Un pool
persistente eliminaría parte de ese sobrecoste sin obligar al usuario a envolver
manualmente el cómputo en una región paralela grande.

La segunda línea es aligerar barreras y reducciones. La espera medida en
\texttt{threading.Condition.wait} llega a porcentajes relevantes con varios
hilos, y cada reducción escalar añade sincronización. Una barrera más ligera en
el runtime nativo y reducciones que eviten barreras completas reducirían el coste
de coordinación que limita CG, FT y MG.

\subsection{Mejoras por benchmark}
Las actuaciones con más probabilidad de impacto no son cambiar únicamente de
versión de Python ni aumentar hilos sin modificar el runtime. La comparación
entre 3.13t, 3.14t y 3.15t muestra curvas muy parecidas cuando la afinidad está
corregida. En CG, la línea útil es agrupar más trabajo entre reducciones y
optimizar barreras. En IS, paralelizar mejor los prefijos acumulados y reducir el
coste de los histogramas por hilo. En MG, adaptar la distribución de trabajo al
tamaño de cada nivel de malla para no activar equipos grandes en niveles con muy
pocas iteraciones útiles. En todos los casos, mantener los bucles críticos
completamente compilables por Cython sigue siendo una condición necesaria.

\subsection{Extensión a BT, SP y LU}
Completar la cobertura con BT, SP y LU permitiría evaluar patrones con
dependencias espaciales más complejas. BT y SP son paralelizables con
\texttt{parallel for} y quedaron fuera por tiempo. LU requiere más que un port
directo. Su barrido SSOR necesita sincronización en \emph{wavefront}, que en
OpenMP se expresa con \texttt{ordered}, \texttt{flush}, \texttt{atomic} o
dependencias entre tareas. Ninguno de esos mecanismos está hoy disponible en las
directivas usadas de OMP4Py (Sección~\ref{sec:benchmarks-adaptados}). Incorporar
alguno de ellos sería una mejora de la biblioteca, no solo una extensión de la
campaña de benchmarks.


\section{Valoración global}
\label{sec:valoracion}

OMP4Py permite expresar paralelismo de estilo OpenMP en Python y producir
ejecuciones correctas sobre benchmarks numéricos no triviales. El resultado más
sólido es la combinación de corrección funcional y mejora por tipado. El
modo~3 reduce los tiempos en más de dos órdenes de magnitud en los casos más
favorables y hace viables cargas que en modo~2 serían demasiado costosas.

El rendimiento paralelo es positivo, pero selectivo. EP y FT muestran speedups
altos con Python~3.15t y rinden parecido en Python~3.13t y Python~3.14t. CG, IS
y MG no escalan bien, ni siquiera con el código tipado y verificado. La utilidad
práctica de OMP4Py no es uniforme. Depende de la combinación entre tipado Cython,
estructura del kernel, configuración de afinidad y eficiencia del runtime.

El valor principal del trabajo es haber portado benchmarks completos, medido sus
caminos de ejecución y separado varios efectos que podrían confundirse,
corrección funcional, mejora por compilación, mejora por tipado, escalabilidad
real, problemas de configuración de entorno y límites propios de cada
benchmark. Esa separación deja una base experimental clara para continuar el
desarrollo de OMP4Py y para evaluar futuras versiones de Python
\emph{free-threaded}. La contribución central no es afirmar que OMP4Py escala
siempre, sino delimitar cuándo escala, cuándo no, y qué parte corresponde al
runtime.
